\documentclass[12pt]{article}
\usepackage{color,soul}
\usepackage{tikz}
\usetikzlibrary{cd}
% Import preambles and macros for notes
\input{../../../../preambles/notes-preamble.tex}
\input{../../../../preambles/notes-macros.tex}
\input{../../../../preambles/theorem-system-section.tex}

\title{MATH 420 Lecture 20}
\author{Deepak Jassal}
\date{March 25\textsuperscript{th}, 2026}

\begin{document}
\maketitle
\section{Exact Sequences}
Exact sequences are of the form
\[
    0\xrightarrow{f_0}A_1\xrightarrow{f_1}\cdots\xrightarrow{f_{k-1}}A_k\xrightarrow{f_k}0
\]
such that Im$(f_j)=\ker(f_{j+1})$
\subsection{Short Exact Sequence}
\[ 
    0\xrightarrow{f_0}A\xrightarrow{f_1}B\xrightarrow{f_2}C\xrightarrow{f_3}0
\]
\ex{
    \[
        0\xrightarrow{f_0}\R\xrightarrow{f_1} C^\infty(\R,\R)\xrightarrow{f_2=\frac{d}{dx}} C^\infty(\R,\R)\to 0
    \]
    $f_1:\R\mapsto$ constant functions. 
}
\ex{
    \[
        0\to\Z\to\Q\to\Q/\Z\to 0
    \]
}
\ex{
    \[
        0\to\Z\to\R\to\R/\Z\cong S^1(\text{sphere})\to0
    \]
}
\ex{
    Selmer group
    \[
        0\to E^{(\Q)}/2E(\Q)\to\text{Sel}_2(E)\to\text{Sha}_E[2]\to0
    \]
    $E$ is an elliptic curve over $\Q$. Computing Sha is an open problem for elliptic curves.
}
\subsection{Commutative Diagrams}
\[
    \begin{tikzcd}[column sep=huge, row sep=huge]
        \mathcal{A} \arrow[r, "F"] \arrow[d, "\alpha"'] & \mathcal{B} \\
        \mathcal{C} \arrow[r, "G"'] & \mathcal{D} \arrow[u, "\beta"']
    \end{tikzcd}
\]
with $F=\beta\circ G\circ\alpha$
\ex{

}
    \[
        \begin{tikzcd}[column sep=large, row sep=large]
            % Row 1: A, B, (empty), (empty)
            X(\Q) \arrow[r, ] \arrow[d, ] & x(\Q_p) \arrow[d, ] & & \\
            % Row 2: C, D, (empty), E
            J(\Q) \arrow[r, ] & J(\Q_p) \arrow[rr, ] & & H^\circ(x_{\Q_p}\Omega')
            \arrow[from=1-2, to=2-4, , bend left=15]
        \end{tikzcd}
    \]
\propp{
    Let $R$ be a ring, and let 
    \[ 
        0\xrightarrow{f_0}M_1\xrightarrow{f_1}M_2\xrightarrow{f_2}M_3\xrightarrow{f_3}0
    \]
    be an exact sequence of $R$ modules. 
    \begin{enumerate}[label=(\alph*)]
        \item If $N\subseteq P$ are submodules of $M_2$ such that $f_1(M_1)\cap N=f_1(M_1)\cap P$ and $f_2(N)=f_2(P)$, then $N=P$;
        \item If $M_1$ and $M_1$ are finitely generaeted, then so it $M_2$;
        \item $M_2$ is Noetherian \textit{iff} $M_1$ and $M_3$ are both Noetherian.
    \end{enumerate}
}{
    \begin{enumerate}[label=(\alph*)]
        \item Let $p\in P$. The second condition implies there exists $n\in N$ such that $f_2(p)=f_2(n)$. Since $f_2$ is a module homomorphism we have $f(n-p)=0$. Therefore, sicne this si an exact sequence, $n-p\in\ker(f_2)\Rightarrow n-p\in$Im$f_1=(f_1(M_1))$. Since $N\subseteq P$ we have $n,\in P$, hence $n-p\in P$. Therefore, $n-p\in f_1(M_1)\cap P$. By the hypothesis $f_1(M_1)\cap P=f_1(M_1)\cap N$. Hence, $n-p\in f_1(M_1)\cap N$, thus $s-n\in N$. But $\underbrace{(p-n)}_{\in N}+\underbrace{n}_{\in N}=p\in N$. Thus, $P\subseteq N$. So $P=N$.
        \item Let $s_1$ be a finite set of generators of $M_1$. Recall, from the property of exact sequences that $f_2:M_2\to M_3$ is surrjective. Since $M_3$ is finitely generated, there exists finite $s_2\subset M_2$ such that $f_2(s_3)$ generates $M_3$. Let $N\subset M_2$ be the submodules generated by $f_1(s_1)$ and $s_3$. Then $f_1(M_1)\cap N=M_2$, and $f_2(N)=M_3=f_2(M_2)$. Part (a) says that $N=M_2$. Hence $M_2$ is generated by $f_1(s_1)\cup s_3$, which is finite.
        \item One direction it clear (i.e., $M_1,M_2$ are Noetherian implies that $M_2$ is Noetherian). Now assume $M_2$ is Noetherian. Let
        \[
            N_1^{(1)}\subseteq N_2^{(1)}\subseteq\cdots
        \]
        be an ascending chain of submodiles of $M_1$. then
        \[
            f_1(N_1^{(1)})\subseteq f_1(N_2^{(1)})\subseteq\cdots 
        \]
        is an ascending chain in $M_2$. Since $M_2$ is Noetherian there exists $k\leq 1$ such that $f_1(N_s^{(1)})=f_1(N_k^{(1)})$ for all $s\geq k$. But $f_1$ is injective, so this implies 
        \[
            N_s^{(1)}=N_k^{(1)}
        \]
        for all $s\geq k$. So $M_1$ is Noetherian. 
    \end{enumerate}
}
\end{document}